\documentclass[addpoints,,12pt]{exam}

\title{Prova.2018.2.Linguagem.Programação.3}

\usepackage[utf8]{inputenc}
\usepackage[brazil]{babel}
\usepackage[T1]{fontenc}
\usepackage{mathtools}
\newcommand*{\renameenviron}[1]{%
  \expandafter\let\csname exam-#1\expandafter\endcsname
      \csname #1\endcsname
  \expandafter\let\csname endexam-#1\expandafter\endcsname
      \csname end#1\endcsname
  \expandafter\let\csname #1\endcsname\relax
  \expandafter\let\csname end#1\endcsname\relax
}
\renameenviron{framed}
\renameenviron{shaded}
\renameenviron{leftbar}
\usepackage{mdframed}
\usepackage{minted}

\mdfsetup{frametitlealignment=\centering}

\usepackage{enumerate}
\usepackage{paralist}
\usepackage{syntax}
\setdefaultleftmargin{0em}{2.2em}{1.87em}{1.7em}{1em}{1em}

\pointpoints{ponto}{pontos}

\newcommand{\figuratopo}{  
	\begin{center}
    	%\vbox {\includegraphics[width=0.20\textwidth]{figuras/logo_topo.jpeg}}
    	\vbox to 0pt{\hfill\fbox{\includegraphics[height=3cm]{figuras/qrcode.png}}}
    	%\includegraphics[scale=0.2]{figuras/qrcode.png}
  	\end{center}}
\renewcommand{\half}{
0.5
}

\begin{document}

%\figuratopo 
\pagestyle{headandfoot}
\firstpageheadrule
\firstpageheader{Prof. Nome do professor}{}{$2^{\text{o}}$ ano do curso - disciplina}
\runningheader{$2^{\text{o}}$ ano do curso - disciplina}{}{}
\runningheadrule
\runningfooter{}{}{p. \thepage\ - \numpages}

\begin{center}
	\bfseries \large Avaliação do $2^{\text{o}}$ Bimestre - 07/06/2018
\end{center}

\begin{mdframed}[userdefinedwidth=1.0\textwidth,align=center,frametitle={Instruções}]
			%\footnotesize
		\begin{enumerate}
			\item Escreva, abaixo, no local indicado o nome completo. 
			\item Esta prova é composta por \numpages\ páginas (incluindo a capa) e \numquestions\ questões.
			\item Aguarde autorização para iniciar a prova. A seguir, antes de iniciar a prova, confira a paginação e o número de questões.
			\item A interpretação das questões é parte do processo de avaliação, não sendo permitidas perguntas.
			\item Não serão permitidos empréstimos de materiais, consultas e comunicação entre os alunos, tampouco o uso de livros e apontamentos. Aparelhos eletrônicos e, em especial, aparelhos celulares deverão ser desligados.
			\item Utilizar folhas em branco para a solução das questões.
			\item As questões podem ser resolvidas e apresentadas fora de ordem, mas devem estar identificadas.
			\item Colocar nomes em todas as folhas utilizadas, inclusive nas que apresentam os enunciados das questões.
			\item Todas as folhas utilizadas devem ser entregues.
			\item Questões resolvidas à lápis podem causar problemas em uma possível revisão de prova. Portanto, prefira respostas a caneta nas partes importantes da solução das questões.
		\end{enumerate}

	\hfill \textbf{Boa Prova!}
	\end{mdframed}
	
	\vfill
	
	\begin{center}
		\textbf{Campo reservado ao professor} \\
		\hqword{Questão}
		\hpword{Pontos}
		\hsword{Nota}
		\gradetable[h]%[pages]  % Use [pages] to have grading table by page instead of question	
	\end{center}
	\vspace{.5cm}
	
	\begin{mdframed}[align=left]
		\textbf{Nome: 12 - NOME COMPLETO DO ALUNO 12}
		\\
		\\Assinatura:
	\end{mdframed}

	\newpage	
	
\noindent \textbf{Questões}

\begin{questions}
\question[10] 

Dada a gramática abaixo, quais das produções passam pelo teste de disjunção?\\
\begin{tabular}{|l c l}
	G: & & \\
	\hline \\
	A & $\rightarrow$ & aB | b | cBB \\
    B & $\rightarrow$ & aB | bA | aBb \\
	C & $\rightarrow$ & aaA | b | caB \\
\end{tabular}
\\
De acordo com cada não terminal da gramática:\\
\\
\begin{choices}
\item A - passa, B - falha, C - passa %verdadeiro 
\item A - falha, B - passa, C - falha
\item A - falha, B - falha, C - falha
\item A - falha, B - passa, C - passa
\item A - passa, B - passa, C - passa
\end{choices}
\answerline


\question[10]
Complete a seguinte frase:\\
_______, 2, 3, 4 e ______. São números inteiros.
\\
\begin{choices}
\item 3, 4 
\item 4, 5 
\item 1, 2 
\item 2, 3
\item 1, 5 %verdadeiro
\end{choices}
\answerline


\question[10]

Os lexemas de uma linguagem de programação não incluem:

\begin{choices}
\item Palavras reservadas
\item Operadores
\item Variáveis %verdadeiro
\item Literais numéricos
\item Identificadores
\end{choices}
\answerline


\question[10]
Analise as frases abaixo e assinale a alternativa que é verdadeira.\\
\\
\begin{choices}
\item TEXTO DA ALTERNATIVA. %verdadeiro
\item TEXTO DA ALTERNATIVA.
\item TEXTO DA ALTERNATIVA.
\item TEXTO DA ALTERNATIVA.
\item TEXTO DA ALTERNATIVA.
\end{choices}
\answerline

\question[10]
Analise as frases abaixo e assinale a alternativa que é verdadeira.\\
\\
\begin{choices}
\item TEXTO DA ALTERNATIVA. %verdadeiro
\item TEXTO DA ALTERNATIVA.
\item TEXTO DA ALTERNATIVA.
\item TEXTO DA ALTERNATIVA.
\item TEXTO DA ALTERNATIVA.
\end{choices}
\answerline

\question[10] 

Dada a gramática abaixo, quais das produções passam pelo teste de disjunção?\\
\begin{tabular}{|l c l}
	G: & & \\
	\hline \\
	A & $\rightarrow$ & aB | b | cBB \\
    B & $\rightarrow$ & aB | bA | aBb \\
	C & $\rightarrow$ & aaA | b | caB \\
\end{tabular}
\\
De acordo com cada não terminal da gramática:\\
\\
\begin{choices}
\item A - passa, B - passa, C - passa
\item A - falha, B - passa, C - passa
\item A - passa, B - falha, C - passa %verdadeiro 
\item A - falha, B - passa, C - falha
\item A - falha, B - falha, C - falha
\end{choices}
\answerline


\question[10]
Analise as frases abaixo e assinale a alternativa que é verdadeira.\\
\\
\begin{choices}
\item TEXTO DA ALTERNATIVA. %verdadeiro
\item TEXTO DA ALTERNATIVA.
\item TEXTO DA ALTERNATIVA.
\item TEXTO DA ALTERNATIVA.
\item TEXTO DA ALTERNATIVA.
\end{choices}
\answerline

\question[10]
Analise as frases abaixo e assinale a alternativa que é verdadeira.\\
\\
\begin{choices}
\item TEXTO DA ALTERNATIVA. %verdadeiro
\item TEXTO DA ALTERNATIVA.
\item TEXTO DA ALTERNATIVA.
\item TEXTO DA ALTERNATIVA.
\item TEXTO DA ALTERNATIVA.
\end{choices}
\answerline

\question[10]
Complete a seguinte frase:\\
_______, 2, 3, 4 e ______. São números inteiros.
\\
\begin{choices}
\item 1, 5 %verdadeiro
\item 1, 2 
\item 2, 3
\item 3, 4 
\item 4, 5 
\end{choices}
\answerline


\question[10]
Complete a seguinte frase:\\
_______, 2, 3, 4 e ______. São números inteiros.
\\
\begin{choices}
\item 1, 5 %verdadeiro
\item 1, 2 
\item 2, 3
\item 3, 4 
\item 4, 5 
\end{choices}
\answerline


\end{questions}

\end{document}
